\section{Conclusion}
Ce stage a permis de livrer une plateforme cohérente et exploitable pour RUHENNA : un espace public orienté valorisation visuelle et prise de contact, et un back-office permettant à l’artisane de gérer son contenu et ses demandes de manière autonome. Le projet répond ainsi à une problématique réelle d’activité indépendante : \textbf{structurer} l’offre, \textbf{centraliser} les demandes, et \textbf{renforcer} la crédibilité en ligne.

\section{Difficultés et leçons apprises}

\subsection{Difficulté principale : apprendre React/Next.js et Prisma en situation réelle}
La difficulté majeure n’a pas été un obstacle isolé, mais une montée en complexité progressive liée au fait de \textbf{ne pas “savoir coder” initialement dans cet écosystème}. En pratique, cela s’est traduit par :
\begin{itemize}
    \item une compréhension progressive du modèle React (composants, état, effets, séparation UI/logique) ;
    \item l’assimilation du paradigme Next.js App Router (layouts, pages, Server Components vs Client Components) ;
    \item l’apprentissage de Prisma comme ORM (schéma, migrations, client typé, contraintes) ;
    \item la mise en relation front/back : concevoir une API, valider les entrées, gérer les statuts, et persister en base.
\end{itemize}

\subsection{Leçons apprises}
Les principales leçons tirées de ce stage sont les suivantes :

\begin{itemize}
    \item \textbf{Anticiper la conception est essentiel} : investir du temps dans la modélisation (schéma de données, architecture) permet de réduire significativement la dette technique et les corrections ultérieures.

    \item \textbf{S’appuyer sur des outils structurants} : l’utilisation de technologies comme le typage statique (TypeScript), les validations d’entrées et les ORM contribue à améliorer la robustesse, la cohérence et la maintenabilité du code.

    \item \textbf{L’importance de la communication entre composants} : une bonne synchronisation entre les différentes couches (front-end, back-end, base de données) est indispensable pour éviter les incompréhensions et les incohérences fonctionnelles.

    \item \textbf{Itérer et tester régulièrement} : valider fréquemment les fonctionnalités permet de détecter rapidement les erreurs et d’ajuster les choix techniques en continu.

\end{itemize}

\section{Perspectives et évolutions possibles}
Le projet est conçu pour évoluer. Les améliorations les plus naturelles sont :
\begin{itemize}
    \item \textbf{E-commerce complet} : panier, commande, gestion stock, frais de livraison.
    \item \textbf{Paiement en ligne} : intégration Stripe (paiements, acomptes pour événements).
    \item \textbf{Comptes utilisateurs} : suivi de commandes, historique de rendez-vous.
    \item \textbf{Emails transactionnels et marketing} : confirmation et rappel de rendez-vous, campagnes promotionnelles.
\end{itemize}

\section{Apports pour l’établissement}
Ce stage a permis de mobiliser et d’illustrer des notions enseignées au cours de la formation, notamment en développement web avec HTML et JavaScript. Il met en évidence l’importance de la structuration des données, de la logique applicative et de la cohérence entre les différentes parties d’un projet web. Il constitue également un exemple concret de mise en application des bonnes pratiques abordées en formation, telles que l’organisation du code, la validation des données et la documentation.

